\begin{proofbox}
	\textbf{\textcolor{CProof}{思路提示}}
	\textbf{证明分两步：首先证明 $y=f(|x|)$ 是偶函数；然后说明其图像由 $y=f(x)$ 在 $x\ge 0$ 的部分通过翻折得到。}

	\medskip
	\textbf{\textcolor{CProof}{正式推导}}
	\begin{description}[style=nextline, leftmargin=2.8em, labelsep=0.5em, itemsep=0.55em]
		\item[\textbf{步骤一（定义域对称性）：}]
			因为 $f(x)$ 的定义域包含 $[0,+\infty)$，所以对任意 $x\in\mathbb{R}$，$|x|\ge 0$ 在定义域内，故 $y=f(|x|)$ 的定义域为 $\mathbb{R}$。
			\[
				D_{f(|x|)} = \mathbb{R}.
			\]
			\par\textbf{理由：} 由条件 $[0,+\infty)\subseteq D_f$；而 $|x|\ge 0$ 恒成立，因此对所有实数 $x$，$|x|\in D_f$，所以 $f(|x|)$ 对所有 $x$ 都有定义。
		\item[\textbf{步骤二（偶函数验证）：}]
			对任意 $x$，计算 $f(|-x|)$。
			\[
				f(|-x|) = f(|x|) \quad (\text{因为 } |-x|=|x|).
			\]
			\par\textbf{理由：} 由偶函数定义 $f(|-x|)=f(|x|)$，所以 $y=f(|x|)$ 是偶函数。
		\item[\textbf{步骤三（分段表达式）：}]
			按自变量符号展开 $f(|x|)$。
			\[
				f(|x|) =
				\begin{cases}
					f(x), & x \ge 0,\
					\[
						2pt] f(-x), & x < 0.
				\end{cases}
				\]
				\par\textbf{理由:} 当 $x\ge 0$ 时 $|x|=x$;当 $x<0$ 时 $|x|=-x$.
			\item[\textbf{步骤四(图像对应关系):}]
				考察图像上的点.
				\begin{itemize}[leftmargin=*]
				\item 当 $x\ge 0$ 时,点 $(x, f(|x|))$ 即 $(x, f(x))$,故 $y=f(x)$ 在 $x\ge 0$ 的部分完全保留.
				\item 当 $x<0$ 时,点 $(x, f(|x|))$ 即 $(x, f(-x))$,与点 $(-x, f(-x))$ 关于 $y$ 轴对称,而 $-x>0$ 属于保留部分.因此 $x<0$ 部分的图像是 $x\ge 0$ 部分图像关于 $y$ 轴的对称复制.
			\end{itemize}
				\par\textbf{理由:} 利用分段表达式,逐段分析点的坐标关系.
		\end{description}

		\medskip
		\textbf{\textcolor{CProof}{结论回扣}}
		\textbf{综上所述,变换 $y=f(|x|)$ 得到偶函数,其图像由原函数 $y=f(x)$ 在 $y$ 轴右侧的部分保留,并关于 $y$ 轴翻折到左侧.}
	\end{proofbox}
